\documentclass{article}
\usepackage{notes}

\setcounter{section}{0}
\renewcommand{\thesection}{4.\arabic{section}}

\begin{document}
\stdtitle{ANALYSIS II}{Exercise Sheet 4}{Jakob Schildhauer}

\section{Derivative computation}
$\p_x u = \sin y \; x^{\sin y - 1}$, $\p_y u = \ln x \; \cos y \; x^{\sin y}$

\section{Connected graphs}
Let $(x_1, y_1), (x_2, y_2) \in \Gamma_f$. 
By connectedness of $U$, there exists $p: [0, 1] \to U, p(0) = x_1, p(1) = x_2, p \in C^0([0, 1])$. 
It follows immediately that $q: t \in [0, 1] \mapsto (p(t), f(p(t))) \in \R^n \times \R^m$ is a member of $C^0([0, 1])$ and $q(0) = x_1, y_1), q(1) = (x_2, y_2)$, so there exists a continuous path between the arbritrary chosen points, proving the graph is connected. 

\section{$p$-norms}
\begin{enumerate}
    \item \dots 
    
    \item $\abs{x}_\infty = \max_{i=1,\ldots,n} \abs{x_i}$
    
    \item \dots 
    
    \item \dots 
    
    \item \dots
    
    \item In this case, $f$ is concave and the triangle inequality fails. 
\end{enumerate}

\section{$p$-means}
\begin{enumerate}
    \item $\mu_{-\infty} = \min_{i=1,\ldots,n} x_i$, $\mu_\infty = \max_{i=1,\ldots,n} x_i$, $\mu_0(x) = 1$
    
    \item \dots 
    
    \item \dots
\end{enumerate}

\section{Mean value for vector-valued functions}
Counterexample: $t \mapsto (\sin \pi t, \sin 2\pi t)$. 

\section{A directional derivative vanish}
\begin{enumerate}
    \item \dots 
    
    \item \dots 
    
    \item \dots
\end{enumerate}
\end{document}
